%%
%% Joule 中文初稿：制造业 AI 的部署位置、企业成本与基础设施足迹
%% 以原 manuscript_cn.tex 为骨架，参考 manuscript.tex 的结果驱动写法。
%% 当前稿为内部初稿；带“待补”的内容不得作为已验证结论引用。
%%

\documentclass[preprint,12pt]{elsarticle}

\makeatletter
\renewcommand\theaffn{\arabic{affn}}
\makeatother

\usepackage[UTF8,fontset=fandol]{ctex}
\usepackage{amssymb,amsmath,amsfonts}
\usepackage{booktabs}
\usepackage{multirow}
\usepackage{tabularx}
\usepackage{makecell}
\usepackage{array}
\usepackage{xcolor}
\usepackage{graphicx}
\usepackage[caption=false,font=scriptsize,labelfont=sf,textfont=sf]{subfig}
\usepackage{url}
\usepackage{xurl}
\usepackage[numbers]{natbib}

\Urlmuskip=0mu plus 1mu
\urlstyle{same}
% 图件由 v0.8.0 workflow 写入 manuscript_figures/；\graphicspath 已指向该目录。
\graphicspath{{../05_results/v0.8.0/result/manuscript_figures/}}

\journal{Joule}

\newcommand{\tbd}[1]{\textcolor{red!70!black}{[待补：#1]}}
\newcommand{\screening}{\textsuperscript{\dag}}

\begin{document}

\begin{frontmatter}

\title{制造业人工智能应在哪里运行？企业部署、云端服务与能源基础设施的联合权衡}

\author{\tbd{作者与单位}}

\begin{abstract}
制造业人工智能（AI）正在从办公辅助扩展到机器视觉、预测维护、生产排程和研发仿真，但未来用电不仅取决于企业采用率，还取决于采用后的任务覆盖、调用强度和计算效率。本文首先综合制造业AI采用和AI能耗证据，再建立一个“相同有效服务”框架，将中国31个制造业大类的六类AI任务映射为CPU/GPU计算、服务器容量、逐时设施负荷和企业成本。我们比较逐厂独立、集团单节点共享和集团多节点协调三种企业内部部署，并以大型集中式云反事实筛查电网、水、土地和建设材料后果。调查证据显示，中国规模以上制造业企业报告使用AI的比例由2023年的20.93\%上升至2026年的30\%以上，但现有资料不足以直接预测唯一的制造业AI用电轨迹；本文因此在v0.8.0中采用透明的低、中、高需求情景。用户成本、代表行业逐时负荷以及电网、水、土地和建设材料结果由对应图件分别呈现。部署位置可能同时改变企业支付和社会基础设施负担；本文不预设本地或云端普遍占优，而是识别这种权衡的来源与证据边界。
\end{abstract}

\begin{keyword}
制造业人工智能 \sep 分布式计算 \sep 云计算 \sep 负荷灵活性 \sep 电网接入容量 \sep 资源足迹
\end{keyword}

\end{frontmatter}

\section*{引言}

制造业AI正在把计算推向生产现场，同时也在放大数据中心的电力、接入、用水、土地和建设压力。因此，问题不只是制造业会用多少电，而是这些计算应在哪里运行、由谁拥有。

集中式云可以池化服务器并提高设施效率，但企业部署还受隐私、时延、可靠性和数据边界约束。若服务器由工厂或集团拥有，任务期限、装机、电费和接入容量可由同一主体协调；这可能带来柔性，却不保证成本更低。比较不同部署方式，必须先固定相同有效服务，再同时看企业付款和社会基础设施后果。

本文对中国31个制造业大类建立计算---能源联合筛查框架，依次回答三个问题：制造业AI需求有多大；逐厂独立、集团单节点共享和集团多节点协调如何改变用户成本；这些选择以及大型集中式云反事实如何重配电网接入、用水、土地和材料足迹。

\section*{结果}

\subsection*{1、制造业AI采用、服务需求与能耗趋势}


\begin{figure}[!t]
\centering
 \includegraphics[width=0.98\textwidth]{figure1_demand_architecture.pdf}
\caption{\textbf{制造业AI需求、异构计算路由与部署架构。}（a）六类制造业AI任务的全国有效服务需求及CPU/GPU结构性路由。（b）完成相同有效服务的逐厂独立（IF）、集团单节点共享（IG-1host）和集团多节点协调（IG-multisite）。（c）v0.8.0低、中、高需求情景的CPU/GPU设施电量构成。（d）31个制造业大类的任务构成和CPU可路由比例。CPU/GPU路由为结构情景而非观测硬件份额；具体数值以v0.8.0已验证图件为准。}
\label{fig:demand_architecture}
\end{figure}

我们首先梳理制造业AI采用、功能覆盖、任务强度和AI能耗的公开证据。中国经济普查资料显示，2023年460,054家规模以上制造业企业中有96,311家报告应用人工智能，对应20.93\%的采用率；行业差异约从11.4\%到36.6\%。工信部在2026年披露，全国规模以上制造业企业人工智能技术应用普及率已超过30\%。这两个时点共同表明AI正在制造业扩散，但统计口径仍需进一步核对，且“应用至少一种AI”既不等于生产环节深度部署，也不直接给出GPU小时、token或电量。

采用后的应用深度同样重要。办公知识、流程Agent、机器视觉、预测维护、生产排程和研发仿真依赖不同的业务驱动量，分别与员工、设备、产线、异常事件和研发作业相关。已有企业调查还显示，采用企业往往先覆盖少数功能，再逐步扩展到多功能和多场址。这意味着未来需求不能由采用企业数乘一个固定AI负荷得到，而应拆成外延采用、任务覆盖、调用强度和单位服务计算效率。通用AI推理效率虽然快速提高，但更高分辨率、更复杂推理和更频繁调用也可能形成反弹效应。

基于这些证据，我们把31个制造业大类映射为六类任务和CPU/GPU计算需求（图~\ref{fig:demand_architecture}a,d），并比较逐厂独立、集团单节点共享和集团多节点协调三种企业内部部署（图~\ref{fig:demand_architecture}b）。v0.8.0保留低、中、高需求情景，用于检验后续用户成本和基础设施结论是否随需求扩大而保持；Figure~1同时给出各情景的设施电量构成。所有需求数值均作为透明情景包络，而不是2030年制造业AI用电的概率预测。



\subsection*{2、负荷角度：部署方式如何改变用户成本}

\begin{figure}[!t]
\centering
 \includegraphics[width=0.98\textwidth]{figure2_enterprise_cost.pdf}
\caption{\textbf{企业内部算力部署的成本与组织边界。}（a）逐厂独立部署（IF）、集团单节点共享池（IG-1host）和集团多节点协调（IG-multisite）的任务与节点边界。（b）三种架构的企业年成本分解。（c）GPU/CPU服务器组数、逐厂离散装机和利用率。（d）集团多节点协调相对集团单节点的增量成本与收益。IF采用逐厂整数装机，两种IG采用连续等效装机，三者完成相同有效AI服务；因此IF与IG的差额包含定容口径，而IG-multisite相对IG-1host主要用于识别跨节点灵活性。}
\label{fig:enterprise_cost}
\end{figure}

用户选择部署方式时，首先比较为相同AI服务承担的付款和自建成本。本地方案包括GPU和CPU服务器及设施年化投资、电量与最大需量费用、运维以及情景中启用的储能；云端方案记录企业购买Token API、GPU容量和CPU容量的付款。由于云端报价已经包含服务商底层服务器、机房和电力成本，这些资源成本不再重复加入企业账单。企业成本和下一节的社会基础设施成本因此是两个不同核算层次。

Figure~2比较中国和美国各自低、中、高制造业AI需求下的本地自建与公开云端采购成本，并分解服务器与设施、电量、最大需量、运维以及云端各类付款。本稿暂不在正文复述全部点估计，而以图件展示成本区间和构成，并据此检验部署选择究竟由资产定容、服务器寿命、云端有效价格还是电力费用驱动。

无论v0.8.0点估计为何，公开目录价都可能高估大型客户的实际云端付款，本地核算也可能遗漏迁移、集成、专业运维、设备残值和技术迭代风险。因此，Figure~2首先用于说明成本构成和决策边界；只有完成服务质量与SLA等价检查和合同价格敏感性后，才解释本地与云端的排序。

用户成本还受到负荷与接入方式影响。Figure~3a以集团单节点共享池作为核心场景，展示代表行业连续168小时的原始生产负荷、AI服务器负荷和新增接入容量。Figure~3b再加入集团多节点协调：集团在多个成员工厂安装算力，并允许符合数据与时延边界的AI任务在截止时间内跨厂调度。两者之差用于度量跨节点灵活性的价值，而不是替换核心场景。

\begin{figure}[!t]
\centering
 \includegraphics[width=0.98\textwidth]{figure3_grid_capacity.png}
\caption{\textbf{集团单节点核心负荷、跨节点调度与生产负荷匹配价值。}（a）集团在一个典型成员工厂安装共享算力池（IG-1host）时，代表行业连续168小时的原始生产负荷、AI服务器负荷、总负荷和新增接入容量。（b）实际负荷下IG-multisite相对IG-1host的跨节点调度价值，并仅对IG-1host以保留原负荷和原负荷置零后分别重新优化的配对差额，分解AI与固定承载工厂生产负荷错峰的成本和容量价值。零负荷是AI独立负荷反事实，不表示工厂停产。}
\label{fig:load_alignment}
\end{figure}

负荷柔性是否能够降低企业购电峰值，还取决于光伏、储能与分时电价的联合运行。代表行业结果表明，集团节点在无AI基准购电、加入AI但尚未经过分布式能源调节的负荷，以及最终购电曲线之间可能出现明显差异；储能充放电、分时电价和荷电状态的联合运行进一步决定峰值是否上升。这一机制用于说明峰值变化的来源，不替代跨行业容量汇总。

\subsection*{3、基础设施角度：部署方式如何改变社会成本与资源足迹}

企业付款并不覆盖部署位置的全部社会后果。相同计算需求集中到大型新建节点时，需要形成新的电网接入、机房壳体、土地和冷却系统；放置在工厂时则可能复用既有建筑和接入设施，同时把冷却用水、室外散热和噪声暴露分散到大量工业节点。由于这些指标尚未采用边界一致的货币化方法，本文把“社会成本”具体呈现为接入容量、现场取水、建筑面积、土地转换、建设材料和潜在噪声影响，而不合成为单一货币总额。

在集团单节点和集团多节点部署下，新增电网接入容量在行业之间高度不均匀。全国等效总量回答各行业对新增容量的贡献，最大代表节点强度则刻画容量集中到单个接入点后的局地压力。两项指标需要并列解读：全国贡献较大的行业不一定具有最高的单节点强度，跨工厂调度也可能改变集团接入峰值。

图~\ref{fig:resource_footprint}比较三种企业内部架构和独立大型云反事实的现场取水与电网接入容量。大型云由全国云情景独立求解，不再以行业单节点代替。建筑、土地和材料账户尚未按新核心架构重建，因此本版不在该图中报告这些量，也不据旧结果陈述其差额。

\begin{figure}[!t]
\centering
 \includegraphics[width=0.98\textwidth]{figure4_resource_footprint.png}
\caption{\textbf{企业内部部署与独立大型云反事实的运行资源需求。}（a）逐厂独立、集团单节点、集团多节点协调及独立全国云反事实的现场取水。（b）对应的新增接入容量。大型云由独立云情景求解，不使用II\_1host代理；材料与土地账户待按新架构重建后另报。}
\label{fig:resource_footprint}
\end{figure}

全国总量之外，影响落在哪里同样重要。工厂侧部署随制造业活动分布到31个省份；云服务情景则按当前智算容量代理重新分配到较少省份（图~\ref{fig:spatial_water}a,c）。两种情景不仅改变省级取水绝对量，也改变取水与地方水稀缺度的叠加。云端情景中，河北、上海、浙江、贵州、内蒙古等省区进入取水量前列，其中较高的AWARE稀缺系数会进一步放大部分地区的稀缺加权影响（图~\ref{fig:spatial_water}b,d）。这一结果说明，集中部署可能把全国平均看似有限的资源需求转化为地方性暴露。

\begin{figure}[!t]
\centering
 \includegraphics[width=0.98\textwidth]{figure5_spatial_concentration.png}
\caption{\textbf{部署位置改变现场取水的空间分布与地方水稀缺压力。}（a,b）工厂侧本地部署的31省取水量及前7省；（c,d）按当前智算容量代理分配的云服务情景。气泡面积和柱长表示估算取水量，颜色表示AWARE2.0年度非农业水稀缺系数。云端省级份额是当前容量空间代理，不是2030规划或云厂商实际流量；当前分析未细分水源和月份。}
\label{fig:spatial_water}
\end{figure}

噪声是同一基础设施边界中尚未量化的维度。大型数据中心把冷却塔、干冷器和备用设备噪声集中在少数园区，分布式机房则可能把较小声源带到更多工厂及其周边。现有项目尚缺可比较的设备声功率、运行占空比、距离衰减和敏感受体数据，因此本稿只将噪声列入社会基础设施成本框架，不报告数值或架构排序。类似地，电网容量是连续MW筛查，土地和材料是设施原型结果，均不等同于已经发生的工程或完整生命周期社会成本。

\section*{讨论}

本文表明，制造业AI“在哪里运行”不能被简化为大型云设施与小型机房的能效比较。即使不同架构提供相同的年度服务，用户成本、单点新增峰值和社会基础设施足迹仍可因节点位置、基础生产负荷和任务期限而显著变化。代表行业的周内负荷结果与全国基础设施筛查共同说明，仅依据年度TWh预测AI基础设施需求会遗漏峰值、空间组织和既有设施复用这些决定性维度。

更重要的是，所有权决定技术柔性能否转化为经济行为。云端客户掌握任务的生产含义和截止期，云服务商控制底层服务器和设施；两者之间的信息和激励并不自动一致。企业自有服务器则把任务安排、设备定容、电费和生产风险置于同一个决策边界。因而，本地部署的主要柔性价值可能不是追逐短时低电价，而是用截止期信息降低服务器装机、复用已有容量并协调生产过程。该机制与传统数据中心需求响应互补：后者强调集中资产在时间和空间上的调度，本文强调任务控制权、资产所有权与电力成本承担者的一致性。

这也改变了对云计算规模经济的理解。集中式云平台确实可能具有更高利用率、更低PUE、更专业的运维和更快的技术更新，但这些优势不必然全部传递为较低的企业付款，也不保证更低的地方接入压力。反过来，企业部署也不能仅凭公开云价较高就被判定为经济。大型客户折扣、网络和存储、系统集成、人员成本、服务器残值、需求不确定性以及模型质量等价性都可能改变打平边界。当前成本结果只能说明“企业自建必然更贵”不是一个稳健先验；最终判断仍需要v0.8.0全国联合成本重算和真实合同边界。

资源足迹进一步提醒我们，分布式并不等于影响消失，而是影响被重新分配。复用既有工厂可以减少绿地建筑和土地转换，却可能把冷却设备、水需求和噪声分散到更多地点；集中设施可能提高冷却效率，却将水、电力和社区影响集中在少数节点。政策和企业评价因此应同时报告全国总量、最大单点、节点数量与当地资源约束，而不是只比较平均PUE或单位计算能耗。对缺水地区、接入受限地区和具有大量可复用工业空间的地区，最优架构可能不同。

总体而言，AI服务器的所有权与位置是连接计算经济性和能源基础设施的关键设计变量。制造业AI不必全部复制大型数据中心的集中扩张路径，也不应被预设为全面回到企业机房。更可行的方向是把部署架构视为一个可优化的组合：在相同服务质量下，根据任务属性、资本利用、能源条件和地方资源约束选择适当层级。该视角将问题从“云端还是本地”推进为“哪些任务在何处运行，谁拥有调度权，以及由此产生的基础设施后果由谁承担”。

\section*{方法}

本文的核算链条从制造业活动与AI采用率出发，依次形成六类AI任务、有效服务需求、CPU/GPU路由、服务器装机和逐时设施负荷（图~\ref{fig:method_workflow}）。后续各方法模块分别定义这一链条的系统边界、服务守恒、容量配置、企业成本和基础设施足迹。

\begin{figure}[!t]
\centering
 \includegraphics[width=0.98\textwidth]{figure1_method.png}
\caption{\textbf{制造业AI需求核算与电力负荷转换方法。}制造业活动和采用率经六类AI任务转化为满足实时性、质量与并发约束的有效服务需求；任务随后路由至CPU或GPU，依据装机裕量与N+1可靠性要求配置服务器，并通过设施附加系数形成逐时设施负荷。}
\label{fig:method_workflow}
\end{figure}

\subsection*{研究边界与部署反事实}

研究对象为中国31个制造业大类的企业AI推理与相关计算服务。功能单位是在共同模型质量、响应时延、吞吐量、可用率和数据边界下完成的有效AI服务。当前主模型不核算基础模型预训练，也不把不同质量的专有模型与开源模型直接视为等价。比较包括：（1）IF，各成员工厂独立安装并仅服务本厂；（2）IG-1host，集团在一个成员工厂设置共享池；（3）IG-multisite，集团在多个成员工厂配置算力并允许符合边界的任务跨厂调度；（4）独立单一绿地云中心，将全国需求置于一个新建大型节点，作为空间集中反事实。云端采购与企业内部物理架构分别建模，避免把采购边界和能源边界混为一谈。

\subsection*{制造业AI有效服务需求}

制造业AI采用和趋势证据优先来自官方统计、企业调查及有明确口径的政策材料。我们分别记录调查对象、参考年份、行业覆盖、采用定义和样本量，并把“使用至少一种AI”“功能覆盖”“生产使用深度”和“员工任务使用”视为不同指标。中国2023年采用率以规模以上制造业企业中报告应用人工智能的单位数除以同期企业总数计算，数据来自《中国经济普查年鉴2023》；2026年“超过30\%”来自工信部制造业数字化转型评估材料，只作为全国扩散锚点，在统计定义完全核对前不与2023年点直接拟合为唯一增长率\citep{NBS2023EconomicCensusYearbook,MIIT2026DigitalTransformationAssessment}。政策目标和市场支出增速仅用于构造情景，不视为实际采用率或电力需求。

AI需求由六类任务构成。对行业$i$、任务$k$和时段$t$，到达的有效服务量写为
\begin{equation}
S^{\mathrm{arr}}_{i,k,t}=D_{i,k}\,a_{i,k}\,q_{i,k}\,p_{i,k,t},
\end{equation}
其中$D_{i,k}$为人员、设备、产线或事件等业务驱动量，$a_{i,k}$为采用外延，$q_{i,k}$为单位驱动量的服务强度，$p_{i,k,t}$为归一化到达分布。低、中、高情景分别改变全国需求尺度和关键任务强度；具体设施电量由v0.8.0全国联合CPU/GPU模型计算。所有情景均保留参数来源、处理步骤和证据等级；IEA和LBNL关于AI及数据中心用电的估计只作数量级校验，不直接作为行业分配数据\citep{IEA2025EnergyAI,Shehabi2024DataCenterEnergy}。

\subsection*{任务时限与调度}

每类任务具有允许执行集合$\mathcal{A}_{k,t}$。实时任务仅能在到达时段执行；批处理或非实时任务可以在截止期内移动。调度量$x_{i,k,t,\tau}$表示到达于$t$的任务在$\tau\in\mathcal{A}_{k,t}$执行，并满足
\begin{equation}
\sum_{\tau\in\mathcal{A}_{k,t}}x_{i,k,t,\tau}=S^{\mathrm{arr}}_{i,k,t}.
\end{equation}
该约束保证所有架构完成相同服务。完全不可调度压力测试令$\mathcal{A}_{k,t}=\{t\}$。这种以截止期约束可延迟计算任务的表示与数据中心时间迁移研究的基本处理一致\citep{Radovanovic2022CarbonAwareComputing}。当前模型采用代表周并按等效周数年化，不表达跨周任务迁移。

\subsection*{异构硬件、容量与三状态设施功率}

任务按照可审计的适用性规则路由到CPU或GPU。对硬件类型$h$，执行服务通过单位服务计算需求转换为加速器或处理器小时。装机容量必须覆盖未平移到达形成的规划下限、优化执行峰值、N+1和统一装机裕量。活动模型以连续等价服务器组表示全国筛查；整数台数和节点最小投资单元作为后续离散敏感性。

设施功率区分关机、在线空闲和计算活动三种状态。时段$t$的IT功率由在线设备空闲功率与随利用率变化的活动功率组成，设施侧功率为
\begin{equation}
P^{\mathrm{fac}}_{n,t}=m^{\mathrm{fac}}_n P^{\mathrm{IT}}_{n,t},
\end{equation}
其中$m^{\mathrm{fac}}_n$为包含冷却和配电损耗的设施倍率，其边界对应PUE中总设施能耗与IT设备能耗之比\citep{DOE2024DataCenterDesign}。该表达允许不同规模设施具有不同效率，同时避免用年度平均PUE替代逐时服务器状态；服务器运行功率和设施效率的范围参考LBNL数据中心能耗核算\citep{Shehabi2024DataCenterEnergy}。

\subsection*{企业负荷、光伏、储能与匹配反事实}

在企业节点$n$，AI设施功率与生产基础负荷、屋顶光伏和储能联合形成净购电：
\begin{equation}
P^{\mathrm{grid}}_{n,t}=P^{\mathrm{base}}_{n,t}+P^{\mathrm{fac}}_{n,t}
-P^{\mathrm{PV}}_{n,t}+P^{\mathrm{ch}}_{n,t}-P^{\mathrm{dis}}_{n,t}.
\end{equation}
储能受到功率、能量、效率和荷电状态约束。制造业生产基础负荷的周内形状来自EWELD工业与商业用户负荷数据集，经行业映射和规模归一化后使用\citep{Liu2023EWELD}。对于每个AI架构，模型在相同电价、屋顶、储能投资和运行规则下求解一个无AI基准。AI增量成本定义为AI情景最优成本减去匹配无AI情景成本。无分布式能源边界把光伏和储能固定为零；有分布式能源边界对两组情景使用相同资源规则，避免将基础负荷本来可以获得的收益计入AI。

\subsection*{企业增量成本}

本地和企业侧架构的年增量成本包括服务器、网络和机房设施年化投资，固定与可变运维，购电电量费用，最大需量费用，以及光伏和储能相对于无AI基准的增量投资与运行费用。资本支出用贴现率$r$和寿命$L$对应的资本回收系数年化：
\begin{equation}
\mathrm{CRF}(r,L)=\frac{r(1+r)^L}{(1+r)^L-1}.
\end{equation}
云端方案记录企业对实例、保留容量、API及必要服务的付款，不再叠加服务商底层服务器、电力或电网成本。当前公开目录价来自各供应商带日期的价格页面，用于透明筛查\citep{AlibabaCloud2026ModelPricing,DeepSeek2026APIPricing,OpenAI2026APIPricing,Anthropic2026ClaudePricing,Google2026GeminiPricing}；大客户折扣、出站流量、存储、迁移、集成和支持通过独立情景处理。除非模型质量、时延和可靠性已验证等价，否则成本比较不被解释为产品优劣。

\subsection*{新增电网接入容量}

节点$n$的既有接入容量被设为匹配无AI基准的最优净购电峰值。新增连续接入容量为
\begin{equation}
\Delta K_n=\max\left\{0,\max_t P^{\mathrm{grid,AI}}_{n,t}
-\max_t P^{\mathrm{grid,noAI}}_{n,t}\right\}.
\end{equation}
全国新增容量为$\sum_n\Delta K_n$，避免容量为云端集中节点与企业侧架构全国总量之差。该物理MW指标不直接进入企业目标函数，也不等同于实际线路或变电站投资。后续工程验证需要加入离散设备档位、项目实现概率、建设时序、电压层级和地区单位造价。

\subsection*{现场水、空间与建设材料}

现场补水以IT电量乘以设施规模和冷却情景对应的水强度估计，水强度采用WUE的现场用水/IT电量边界\citep{Mytton2021DataCentreWater}。当前指标不包含发电侧用水，也尚未在所有情景中区分取水与耗水。省级稀缺加权使用AWARE2.0年度非农业表征因子\citep{Seitfudem2025AWARE2}。空间核算区分设备占用面积、新增建筑面积和新增土地；既有厂房复用、园区扩建与绿地新建被视为互斥实现情景。建设材料仅对新增结构壳体核算混凝土、钢筋和结构钢，小型机房改造材料在缺乏样本时记为未报告，而非假定为零。水、土地和材料分别报告，不进行证据不足的货币化加总。

\subsection*{敏感性、验证与推断边界}

敏感性分析首先对需求规模、单位服务算力、任务柔性、CPU/GPU性能与功率、设施倍率、装机裕量、可靠性、服务器寿命、云端有效价格、电价和最大需量价格进行结构化低---中---高或单因素筛查。当前未获得可靠概率分布的参数不进入概率性Monte Carlo或Sobol解释。模型验证包括服务守恒、容量可行性、逐时能量平衡、无AI反事实匹配、成本分项重构、31行业与架构覆盖以及全国汇总守恒。数值结果始终与模型版本、配置和输入快照共同引用。

本文的推断限于情景比较和机制筛查。需求结果不是统计预测；连续MW不是具体电网工程；公开云价格不是企业平均合同；单一云节点不是服务商实际网络；资源足迹MVP也不是完整生命周期评价。只有在相同服务、边界一致且输入可追溯时，架构之间的差异才被解释为部署位置或所有权的结果。

\subsection*{Limitations}

本文仍有五类重要局限。第一，制造业AI需求是基于业务驱动量和外部总量校验的情景构造，不是对未来企业采用的统计预测。第二，代表周、合成或跨地区借用的行业负荷曲线不能完全描述季节性、生产停机和地区差异。第三，连续等价服务器组低估了小型节点的整数装机、冗余和运维门槛。第四，大型集中节点和云端空间分配是用于识别集中度效应的反事实，不代表真实云服务商的多区域网络。第五，成本、水、材料、电网容量和噪声处于不同成熟度，本文没有证据将其货币化或声明综合社会成本最优。全文只接受v0.8.0结果；成本与下游足迹必须在同一依赖链完成后整体更新，不能混入旧版本或局部替换。

尽管如此，该框架提供了一条更符合企业和能源系统共同决策的研究路径。未来工作应优先获得企业AI使用轨迹、任务期限、CPU/GPU性能、真实云合同和本地运维案例；在能源侧加入全年多季节负荷、离散配变与馈线投资、云端候选节点、水源和流域季节压力。对于企业，部署决策可采用分层架构：隐私和低时延任务留在工厂，可共享的稳定任务由集团协调，需求尖峰和专有模型服务使用公有云。对于电网规划者，新增AI需求评估则应从“多少TWh”扩展为“什么任务、由谁调度、落在哪些节点、形成多大峰值”。

\section*{数据与代码可获得性}

制造业任务参数、处理后的服务量、活动配置、模型输出和图件生成脚本保存在项目的数据、配置、结果与代码目录中。投稿前将提供可公开部分的永久仓库链接、版本号和许可证。受许可或商业保密限制的数据将提供来源、处理方法和可复现的替代输入。\tbd{公开仓库、DOI和数据许可声明。}

\section*{致谢}
\tbd{资助、合作与数据支持。}

\section*{作者贡献}
\tbd{按CRediT taxonomy补充。}

\section*{利益冲突声明}
作者声明不存在利益冲突。\tbd{由全体作者确认。}

\bibliographystyle{elsarticle-num-names}
\bibliography{reference}

\end{document}
